%-------------------------------------------------------------------------------
\section{Architecture Overview}
\label{sec:methodology}
\label{sec:method-overview}
%-------------------------------------------------------------------------------

Web PKI normative sources serve as the input to the framework. Through keyword recall and constrained extraction, candidate rules are converted into a structured representation. As shown in Fig.~\ref{fig:architecture}, the framework contains four modules.

First, the knowledge graph construction module handles scattered definitions, inherited requirements, and cross-document references by organizing the collected sources into a traceable source context for rule extraction.

Second, the deterministic context retrieval module prevents the LLM from inventing unsupported context by supplying source-grounded context for each target section in a reproducible way.

Third, DLEP balances completeness and precision by first recalling candidate normative rules; for precise extraction, the second layer uses an LLM to convert them into a structured representation.

Fourth, the four-condition lintability framework determines whether a rule can be implemented as a static lint check.

\begin{figure*}[tb]
  \centering
  \paperfig{figures/fig01_architecture.png}{Overall architecture of the framework, showing knowledge-graph construction, deterministic
    context retrieval, dual-layer rule extraction, and the four-condition
    lintability framework.}
  \caption{Overall Architecture of the Framework.}
  \label{fig:architecture}
\end{figure*}
